% Answer key for the 2.1-2.5 Study Guide -- August 2026
% Numbering matches the study guide (1-32, then Tommy 1 and Tommy 2).
% Textbook answers only, no work shown; Tommy's questions are worked in full.
%
% PRIVATE MAP (comments never appear in the PDF; delete this block if you ever
% hand out the .tex source itself).  Key number = textbook problem:
%   1=2.1#1   2=2.1#5   3=2.1#9   4=2.1#12  5=2.1#13  6=2.1#14  7=2.1#16
%   8=2.1#17  9=2.1#21  10=2.1#25 11=2.1#33 12=2.1#35 13=2.1#37
%   14=2.3#13 15=2.3#45 16=2.3#49 17=2.3#62
%   18=2.4#7  19=2.4#9  20=2.4#21 21=2.4#23 22=2.4#39
%   23=2.5#1  24=2.5#3  25=2.5#7  26=2.5#8  27=2.5#10
%   28=2.5#20 29=2.5#23 30=2.5#27 31=2.5#33 32=2.5#55
%
\documentclass[11pt]{article}
\usepackage[margin=1in]{geometry}
\usepackage{amsmath,amssymb}
\usepackage{enumitem}
\usepackage{multicol}

\newlist{ans}{itemize}{1}
\setlist[ans]{leftmargin=3.2em,itemsep=4pt,topsep=4pt,label={}}
\newcommand{\ansitem}[1]{\item[\textbf{#1.}]}

\title{\textbf{2.1--2.5 Study Guide}\\[2pt] \textbf{Answer Key}}
\author{}
\date{August 2026}

\begin{document}
\maketitle
\vspace{-2em}

\begin{multicols}{2}
\begin{ans}
  \ansitem{1} $-7$
  \ansitem{2} $7$
  \ansitem{3} $-3$
  \ansitem{4} $4$
  \ansitem{5} $4$
  \ansitem{6} $19$
  \ansitem{7} $-4$
  \ansitem{8} $32$
  \ansitem{9} $12$
  \ansitem{10} (a) $-1$\quad (b) $1$\quad (c) DNE
\end{ans}
\end{multicols}

\begin{ans}
  \ansitem{11} (a) $1$\quad (b) $1$\quad (c) $1$\quad (d) $3$\quad (e) $3$\quad (f) $3$
  \ansitem{12} (a) $1$\quad (b) $0$\quad (c) DNE\quad (d) $1$\quad (e) $0$\quad (f) DNE
  \ansitem{13} (a) DNE $(f(x)\to-\infty)$\quad (b) DNE $(f(x)\to\infty)$\quad (c) DNE\quad
                (d) DNE (nothing is graphed to the left of $x=0$)\quad (e) $0$\quad (f) DNE
\end{ans}

\begin{multicols}{2}
\begin{ans}
  \ansitem{14} $-13$
  \ansitem{15} $1$
  \ansitem{16} (a) $0$\quad (b) DNE\quad (c) DNE
  \ansitem{17} $0$
  \ansitem{18} (a) $\infty$\quad (b) $-\infty$\quad (c) DNE
  \ansitem{19} (a) $\infty$\quad (b) $\infty$\quad (c) $\infty$
  \ansitem{20} $1$
  \ansitem{21} DNE
\end{ans}
\end{multicols}

\begin{ans}
  \ansitem{22} Sketch: horizontal asymptote $y=-2$ at both ends; vertical asymptotes
        $x=-1$ and $x=3$; the graph falls to $-\infty$ on the left of $x=-1$, rises to
        $\infty$ on the right of $x=-1$, rises to $\infty$ on the left of $x=3$, and
        falls to $-\infty$ on the right of $x=3$.
\end{ans}

\begin{ans}
  \ansitem{23} Jump discontinuity at $x=2$
  \ansitem{24} Removable discontinuity at $x=2$
  \ansitem{25} Infinite discontinuity at $x=2$
  \ansitem{26} Infinite discontinuity at $x=2$
  \ansitem{27} Removable discontinuity at $x=0$
  \ansitem{28} Continuous: $f(-5)=3$ and $\displaystyle\lim_{x\to-5}f(x)=3$
  \ansitem{29} $f(-2)$ is undefined
  \ansitem{30} $\displaystyle\lim_{x\to3}f(x)=1$ but $f(3)=0$
  \ansitem{31} $x=-2$ and $x=1$
  \ansitem{32} $f(-1)=0$, $f(2)=9$; for any $w$ with $0\le w\le9$, take $c=\sqrt[3]{w-1}$,
        which lies in $[-1,2]$
\end{ans}

%======================================================================
\section*{Tommy's Questions --- Full Solutions}
\noindent\textit{These two are ours, so the steps are shown.}
\medskip

\subsection*{Tommy 1}

\noindent\textbf{(a)} Both top and bottom are $0$ at $x=3$, so factor first:
\[
\lim_{x\to3}\frac{x^{2}-x-6}{x^{2}-9}
=\lim_{x\to3}\frac{(x-3)(x+2)}{(x-3)(x+3)}
=\lim_{x\to3}\frac{x+2}{x+3}
=\frac{5}{6}.
\]
Cancelling $x-3$ is legal because $x\to3$ means $x\neq3$, so $x-3\neq0$.

\medskip
\noindent\textbf{(b)} Multiply by the conjugate:
\[
\frac{\sqrt{9+h}-3}{h}\cdot\frac{\sqrt{9+h}+3}{\sqrt{9+h}+3}
=\frac{(9+h)-9}{h\left(\sqrt{9+h}+3\right)}
=\frac{h}{h\left(\sqrt{9+h}+3\right)}
=\frac{1}{\sqrt{9+h}+3}.
\]
\[
\lim_{h\to0}\frac{\sqrt{9+h}-3}{h}=\frac{1}{\sqrt{9}+3}=\frac{1}{6}.
\]

\medskip
\noindent\textbf{(c)} Write $2x-6=2(x-3)$ and split the absolute value.

\noindent (i) For $x<3$: $|x-3|=-(x-3)$, so
$f(x)=\dfrac{2(x-3)}{-(x-3)}=-2$, hence $\displaystyle\lim_{x\to3^{-}}f(x)=-2$.

\noindent (ii) For $x>3$: $|x-3|=x-3$, so
$f(x)=\dfrac{2(x-3)}{x-3}=2$, hence $\displaystyle\lim_{x\to3^{+}}f(x)=2$.

\noindent (iii) The one-sided limits are $-2$ and $2$, which are not equal, so
$\displaystyle\lim_{x\to3}f(x)$ \textbf{does not exist}.

\medskip
\noindent\textbf{(d)} For every $x\neq3$, $\left|\sin\!\frac{1}{x-3}\right|\le1$, so
\[
-(x-3)^{2}\;\le\;(x-3)^{2}\sin\!\frac{1}{x-3}\;\le\;(x-3)^{2}.
\]
Since $\displaystyle\lim_{x\to3}-(x-3)^{2}=0$ and $\displaystyle\lim_{x\to3}(x-3)^{2}=0$,
the sandwich theorem gives
$\displaystyle\lim_{x\to3}(x-3)^{2}\sin\!\frac{1}{x-3}=0$.

\bigskip
\subsection*{Tommy 2}

\noindent Factor once and use it for every part:
\[
f(x)=\frac{x^{2}-4}{x^{2}-x-6}=\frac{(x-2)(x+2)}{(x-3)(x+2)}=\frac{x-2}{x-3}
\qquad\text{for }x\neq-2,\ x\neq3 .
\]

\medskip
\noindent\textbf{(a)} The denominator $x^{2}-x-6$ is $0$ at $x=-2$ and $x=3$, so those are the
only numbers where $f$ is discontinuous.
\[
\lim_{x\to-2}f(x)=\lim_{x\to-2}\frac{x-2}{x-3}=\frac{-4}{-5}=\frac{4}{5},
\]
a finite limit, so $x=-2$ is a \textbf{removable} discontinuity (a hole at $\left(-2,\tfrac45\right)$).
At $x=3$ the factor $x-3$ does not cancel and $f(x)\to\pm\infty$ (see part (b)), so $x=3$ is an
\textbf{infinite} discontinuity.

\medskip
\noindent\textbf{(b)} Near $x=3$ the numerator $x-2\to1>0$.
\[
x\to3^{-}:\ x-3\to0^{-}\ \Rightarrow\ \lim_{x\to3^{-}}f(x)=-\infty,
\qquad
x\to3^{+}:\ x-3\to0^{+}\ \Rightarrow\ \lim_{x\to3^{+}}f(x)=\infty.
\]

\medskip
\noindent\textbf{(c)} Divide numerator and denominator by $x$:
\[
\lim_{x\to\pm\infty}\frac{x-2}{x-3}
=\lim_{x\to\pm\infty}\frac{1-\dfrac{2}{x}}{1-\dfrac{3}{x}}
=\frac{1-0}{1-0}=1 .
\]
So $\displaystyle\lim_{x\to\infty}f(x)=\lim_{x\to-\infty}f(x)=1$.
Horizontal asymptote: $y=1$. Vertical asymptote: $x=3$ \textbf{only} --- $x=-2$ is a hole,
not an asymptote, because the limit there is finite.

\medskip
\noindent\textbf{(d)} $f$ is continuous on $[0,2]$ (its only discontinuities, $-2$ and $3$, are
outside that interval), and
\[
f(0)=\frac{0-2}{0-3}=\frac{2}{3},\qquad f(2)=\frac{2-2}{2-3}=0 .
\]
Since $0\le\tfrac12\le\tfrac23$, that is $f(2)\le\tfrac12\le f(0)$, the intermediate value
theorem guarantees a $c$ in $[0,2]$ with $f(c)=\tfrac12$. Solving:
\[
\frac{c-2}{c-3}=\frac{1}{2}
\ \Longrightarrow\ 2(c-2)=c-3
\ \Longrightarrow\ 2c-4=c-3
\ \Longrightarrow\ c=1,
\]
and $c=1$ does lie in $[0,2]$. Check: $f(1)=\dfrac{1-2}{1-3}=\dfrac{-1}{-2}=\dfrac12$. \checkmark

\end{document}
